\section{The Function Generator's Control Network}
\label{sec:control-network}

The function generator is the only source whose output is not simply a fixed preset: its
amplitude and frequency are set by two small, undirected utility blocks, a voltage module
and a frequency module, rather than by right-clicking the generator itself (which instead
only cycles its waveform shape -- sine, square, or triangle -- a purely local preset that
never propagates). This section describes the mechanism by which a module's setting reaches
one or more generators, because it is easy to mistake for an extension of the electrical
topology described in Section~\ref{sec:architecture} when it is, in fact, a second and
entirely independent network built on the same primitive (block-to-block adjacency) for a
different purpose.

\subsection{Why a second network}

Section~\ref{sec:limitations} notes that every adjustable quantity in the mod, amplitude and
frequency included, is one of a small number of discrete presets rather than a numeric
entry, to avoid introducing a text-input interface foreign to the base game. A single
right-click-to-cycle control on the generator itself would satisfy that constraint for one
generator, but a classroom exercise frequently calls for several generators sharing one
setting -- for instance, driving multiple parallel RC networks with an identical input to
compare their responses side by side, an exercise Section~\ref{sec:pedagogy} discusses.
Requiring every generator in such a bank to be clicked to the same preset individually, and
kept in sync by hand as the exercise evolves, would be tedious and error-prone. The voltage
and frequency modules solve this by letting a single click on one module drive every
generator reachable through a chain of like-purpose blocks, the same one-to-many relationship
a bench power supply's daisy-chained ``sync'' output provides to several instruments at once.

\subsection{Modules: undirected cubes, not electrical components}

A voltage module and a frequency module are both instances of the same abstract block
entity, differing only in which of two channels, voltage or frequency, they \emph{author}.
Neither has an orientation or electrical leads: unlike a resistor or diode, which expose
exactly two faces as leads and insulate the remaining four, a module treats all six of its
faces identically, and a module may attach to a generator, or to another module, on any
face -- including a generator's own electrical lead faces, since this contact is never
folded into the union--find partitioning that Section~\ref{sec:architecture} describes and
never contributes a term to the MNA system in Section~\ref{sec:circuit-solver}. A module's
own adjustable quantity is a preset cycled by an empty-handed right-click, exactly as a
resistor's resistance or a source's waveform is: a voltage module cycles among
$\{1.5, 5, 9, 12, 24\}$~V and a frequency module among $\{0.5, 1, 2, 5, 10\}$~Hz.

\subsection{Authorship and one-hop-per-tick relay}

Every module, regardless of which channel it authors, stores \emph{both} channels as a pair
$(\text{value}, \tau)$, where $\tau$ is the game tick at which that channel's value was most
recently \emph{authored} -- either directly, when a player cycles that module's own preset
(which always re-stamps $\tau$ to the current tick, even if the cycle happens to land back
on the same value), or indirectly, by adopting a neighboring module's more recently authored
pair. Once per game tick, independently for each channel, every module compares its own
stored $\tau$ against every immediately adjacent module's $\tau$ for that same channel and
adopts whichever pair carries the strictly larger tick:
\begin{equation}
(\text{value}, \tau) \leftarrow \operatorname*{arg\,max}_{\tau' \,\in\, \{\tau\} \,\cup\,
\{\tau_n : n \in \mathcal{N}\}} \tau',
\end{equation}
where $\mathcal{N}$ is the set of the module's face-adjacent neighboring modules. A tie is
resolved in favor of the module's own current value (no oscillation results from two
equally recent authors reaching a module simultaneously), and a module with no neighbor at
all simply keeps relaying its last-known pair indefinitely, including after a neighbor is
removed -- there is no explicit ``disconnected'' state, only a value that stops receiving
fresher ticks. After updating both of its channels, a module unconditionally pushes its
resolved voltage and frequency to every function generator directly adjacent to it.

Because each module consults only its immediate neighbors' \emph{already-updated} values
from the previous tick, and does not look further ahead along the chain, a freshly
authored value advances exactly one block per game tick, a discrete diffusion rather than
an instantaneous update. At 20 ticks per second this delay is imperceptible for any chain a
player is likely to build by hand (a twenty-block chain, for instance, fully converges in
one second), but it is a real, measurable propagation time rather than an idealization, and
is mentioned here specifically because it is the one place in the mod where an in-world
distance between two blocks has a temporal consequence, unlike the electrical network, whose
solve is global and instantaneous within a tick regardless of how large the constructed
circuit is.

\subsection{Mixed chains and multiple generators}

A single physical chain may freely mix voltage and frequency modules: each block relays both
channels regardless of which one it authors, so a chain alternating between the two kinds
still carries a single, globally agreed voltage value and a single, globally agreed
frequency value end to end, authored independently by whichever module of each kind was most
recently clicked anywhere in the chain. Any number of generators may attach to any module in
the chain, at any point along it, and each receives the same converged amplitude and
frequency every tick; the push to a generator is idempotent (a generator only remarks its
network dirty, prompting a resynchronization to clients, if the incoming value actually
differs from its current one), so a large bank of generators sharing one control imposes no
per-tick cost beyond the assignment itself. A generator with no module attached at all
retains its default amplitude and frequency (5~V, 1~Hz) rather than becoming unset, so a
newly placed generator is immediately usable before any module is ever placed.
